\chapter{Valutazione Economica per gli Investitori}
\label{cap:valutazione}

% DA COMPLETARE: l'intero capitolo dipende dai risultati della simulazione (Capitolo 4),
% che non sono ancora disponibili. La struttura e' predisposta.

Questo capitolo traduce i risultati della simulazione in grandezze economiche, dal punto di
vista di un investitore che debba decidere se e quanto capitale impegnare in un impianto di
accumulo.

\section{Metriche finanziarie}
\label{sec:metriche}

% DA COMPLETARE: definizione delle metriche.
% Elementi da trattare:
%   - ricavo lordo annuo da arbitraggio, ottenuto sommando sui periodi il margine
%     calcolato con i prezzi controfattuali della simulazione;
%   - costi: investimento iniziale (CAPEX, in euro per MWh installato), costi operativi
%     annui (OPEX), costo di degrado gia' introdotto nella Sezione 3.3;
%   - indicatori: valore attuale netto (VAN), tasso interno di rendimento (TIR),
%     tempo di ritorno dell'investimento, e LCOS (Levelized Cost of Storage);
%   - orizzonte temporale e tasso di sconto da adottare, con giustificazione.

\section{Risultati di redditività}
\label{sec:redditivita}

% DA COMPLETARE: risultati, da produrre dopo la simulazione.
% Elementi attesi:
%   - redditivita' per ciascuno scenario di potenza e durata definito nella Sezione 4.3;
%   - individuazione del dimensionamento ottimale, cioe' della combinazione di potenza e
%     durata che massimizza il valore attuale netto;
%   - punto centrale della tesi: la redditivita' marginale decresce al crescere della
%     capacita' installata, perche' l'accumulo comprime da se' il differenziale di prezzo
%     da cui trae profitto. Da quantificare la capacita' oltre la quale l'investimento
%     marginale non e' piu' remunerativo.

\section{Analisi di sensitività}
\label{sec:sensitivita}

% DA COMPLETARE: analisi di sensitivita' dei risultati alle ipotesi.
% Parametri rispetto a cui far variare i risultati:
%   - rendimento di ciclo completo;
%   - costo di investimento per MWh;
%   - tasso di sconto;
%   - costo di degrado per ciclo;
%   - anno di riferimento dei prezzi, per distinguere l'effetto delle condizioni di
%     mercato di un anno particolare da quello strutturale;
%   - assunzioni del modello di mercato: in particolare il perimetro zonale (Sezione 4.1)
%     e l'ipotesi che i flussi di scambio con l'esterno non reagiscano al prezzo.
